% ID: SEL-D0059-OB-E01
% ESTADO: cotejada. Segunda pasada visual y borrador compilado revisados.
% Fuente: Selectividad/sel_2011_matematicas_aplicadas/222-EXAMEN 2.pdf
% SHA-256: 998f20402054ca6d0454da83ffc8e1ff6951a576b48836a32e81e524c5985214
% Página: 2; opción B; ejercicio original 1.

\begin{enumerate}[label=\alph*)]
\item \textbf{(1.5 puntos)} Dibuje el recinto del plano definido por el siguiente sistema de inecuaciones y determine sus vértices:
\[
y\geq200-2x,\quad x-100\leq3y,\quad x+2y\leq600,\quad x\geq0.
\]
\item \textbf{(1 punto)} Sabiendo que $A(0,2)$, $B(1,4)$, $C(3,4)$, $D(4,2)$ y $E(2,1)$ son los vértices de una región factible, determine en ella el mínimo y el máximo de la función $F(x,y)=10x+5y+21$, e indique los puntos donde se alcanzan.
\end{enumerate}
